
\documentclass{article}
\usepackage{colortbl}
\usepackage{makecell}
\usepackage{multirow}
\usepackage{supertabular}

\begin{document}

\newcounter{utterance}

\centering \large Interaction Transcript for game `ifeval', experiment `ifeval', episode 158 with llama{-}3.1{-}8b{-}instruct.
\vspace{24pt}

{ \footnotesize  \setcounter{utterance}{1}
\setlength{\tabcolsep}{0pt}
\begin{supertabular}{c@{$\;$}|p{.15\linewidth}@{}p{.15\linewidth}p{.15\linewidth}p{.15\linewidth}p{.15\linewidth}p{.15\linewidth}}
    \# & \multicolumn{2}{c}{Player} && \multicolumn{2}{c}{Game Master} \\
    \hline

    \theutterance \stepcounter{utterance}  
    & & & \multicolumn{4}{p{0.6\linewidth}}{
        \cellcolor[rgb]{0.9,0.9,0.9}{
            \makecell[{{p{\linewidth}}}]{
                \texttt{\tiny{[P1$\langle$GM]}}
                \texttt{I was hoping you could help me out with a few things. I wish to learn about how to pay my taxes. Can you summarize the process for me? Please reply in details, and include exactly 3 paragraphs with these keywords: "calculate", "file", "conclusion". Separate the paragraphs with ***.} \\
            }
        }
    }
     \\ \\

    \theutterance \stepcounter{utterance}  
    & \multicolumn{4}{p{0.6\linewidth}}{
        \cellcolor[rgb]{0.561,0.737,0.561}{
            \makecell[{{p{\linewidth}}}]{
                \texttt{\tiny{[P1$\rangle$GM]}}
                \texttt{Here's a detailed summary of the tax payment process:} \\
\\ 
\texttt{To pay your taxes, you'll first need to calculate your tax liability. This involves determining the amount of income tax you owe to the government. You can do this by using tax software or consulting with a tax professional. You'll need to gather all necessary documents, including your W{-}2 forms, 1099 forms, and any other relevant tax{-}related documents. You'll also need to determine your filing status, which can affect your tax rate and deductions. Once you have all the necessary information, you can use tax software or a tax calculator to determine your tax liability.} \\
\\ 
\texttt{*** } \\
\\ 
\texttt{Next, you'll need to file your tax return with the government. This typically involves submitting a tax return form, such as Form 1040, along with any supporting documentation. You can file your tax return electronically or by mail, depending on your preference. If you're filing electronically, you can use tax software or the IRS's e{-}file system. If you're filing by mail, you'll need to print out your tax return form and supporting documentation, and mail it to the IRS. Be sure to keep a copy of your tax return for your records, as you may need to refer to it in the future. The deadline for filing your tax return is typically April 15th of each year, although this can vary depending on your filing status and other factors.} \\
\\ 
\texttt{*** } \\
\\ 
\texttt{Once you've filed your tax return, you'll need to pay any tax liability you owe. You can do this by making a payment online, by phone, or by mail. You can also set up a payment plan with the IRS if you're unable to pay your tax liability in full. If you've overpaid your taxes, you may be eligible for a refund. In this case, you can file for a refund and receive your money back. In conclusion, paying your taxes involves calculating your tax liability, filing your tax return, and paying any tax liability you owe. By following these steps, you can ensure that you're in compliance with tax laws and regulations, and avoid any potential penalties or fines.} \\
            }
        }
    }
    & & \\ \\

    \theutterance \stepcounter{utterance}  
    & & & \multicolumn{2}{p{0.3\linewidth}}{
        \cellcolor[rgb]{0.9,0.9,0.9}{
            \makecell[{{p{\linewidth}}}]{
                \texttt{\tiny{[GM$|$GM]}}
                {'keywords:existence': {'keywords': ['calculate', 'file', 'conclusion']}, 'length_constraints:number_paragraphs': {'num_paragraphs': 3}}
            }
        }
    }
    & & \\ \\

    \theutterance \stepcounter{utterance}  
    & & & \multicolumn{2}{p{0.3\linewidth}}{
        \cellcolor[rgb]{0.9,0.9,0.9}{
            \makecell[{{p{\linewidth}}}]{
                \texttt{\tiny{[GM$|$GM]}}
                \texttt{game\_result = WIN} \\
            }
        }
    }
    & & \\ \\

\end{supertabular}
}

\end{document}
